\documentclass[10pt,a4paper]{amsart}
\usepackage[margin=19mm]{geometry}
\usepackage{amsmath,amssymb,mathtools}
\usepackage[scr=rsfs,cal=euler]{mathalfa}
\usepackage{xfrac}
\usepackage[unicode,hidelinks,bookmarks=false]{hyperref}
\newcommand{\ie}{i.e.\ }
\newcommand{\cf}{cf.\ }
\newcommand{\resp}{respectively\ }
\newcommand{\Subsection}{\subsection}
\providecommand{\nobreakdash}{\nobreak\hbox{-}}
\setlength{\emergencystretch}{2em}
\multlinegap0pt
\usepackage{bm}
\usepackage{bbm}
\usepackage{dsfont}
\usepackage{xspace}
\usepackage{cancel}
\usepackage{mathtools}
\usepackage{amsthm}
\makeatletter
\@ifundefined{Th}{\newtheorem{Th}{Th}}{}
\@ifundefined{Lem}{\newtheorem{Lem}[Th]{Lemma}}{}
\makeatother
\title[Large Coupling Convergence Beyond Definiteness]{Large Coupling Convergence Beyond Definiteness}
\author{Christian Koke}
\date{Source: arXiv:2601.18055v2 (CC BY 4.0)}
\begin{document}
\maketitle

\noindent\emph{Source and licence.} Large Coupling Convergence Beyond Definiteness. Version arXiv:2601.18055v2 (CC BY 4.0), distributed under the Creative Commons Attribution 4.0 International licence, as recorded in the arXiv licence metadata for this exact version and shown on the abstract page https://arxiv.org/abs/2601.18055v2. Full rights notice: licenses/arxiv-2601.18055-CC-BY-4.0.txt.\par\smallskip

\noindent\textit{Editorial scope.} Counted unit: the Lem whose source label is \texttt{conv\_to\_projection} in the article (source0.tex lines 724--730, source label \texttt{conv\_to\_projection}), delivered together with the article's own complete original proof of it (source0.tex lines 731--798, one \texttt{proof} environment of 68 lines). No mathematics is written, repaired or re-derived: statement and proof are the article's own text line for line. Required context retained uncounted: none. Every definition, display and label of those lines is retained unchanged. Omitted from this package and not redistributed: 1--723, 799--1846. Nothing inside a retained excerpt is omitted or altered: every retained body is delivered exactly as the article's lines are written, so no omission receipt of any kind accompanies this package. Citations: the retained excerpts carry no citation command at all, so no bibliography entry of the article is transcribed and no reference is redistributed. Reference adaptation: none was needed and none was made; every reference inside the delivered unit resolves inside it, so no unresolved reference or citation is produced. Printed slips kept literal: no printed word, symbol, hypothesis or number of the counted statement or of its proof was altered. Layout and class: the article's own class is \texttt{amsart} with packages including amsmath, amssymb, amsthm, graphicx, mathrsfs, url, color, hyperref, floatrow, xcolor, tikz, slashed, amsfonts, calrsfs; the wrapper is the lane's pinned \texttt{amsart} editorial wrapper at 10pt on A4, which restates the theorem-like environments the retained text needs and adds the article's own macro definitions that the retained text needs (none), with extra wrapper packages (bm, bbm, dsfont, xspace, cancel, mathtools). The delivered numbering is the unit's own. Assets: no retained span contains a figure environment, a graphics-inclusion command, a PDF-page inclusion, a table or a TikZ picture; no image file is copied into this package, the package declares no asset dependency and no image file is redistributed with it. Licence: version arXiv:2601.18055v2 of this article is distributed under the Creative Commons Attribution 4.0 International licence, as recorded in the arXiv licence metadata for this exact version and shown on the abstract page https://arxiv.org/abs/2601.18055v2. A full rights notice is delivered at licenses/arxiv-2601.18055-CC-BY-4.0.txt. Characters: every retained excerpt is pure ASCII, so the delivered engine, XeLaTeX with its default Latin Modern text font, reports no missing character.
\begin{Lem}\label{conv_to_projection}
	Let $B: \mathcal{D}(B) \rightarrow \mathcal{H}$ be closed  with  isolated eigenvalue $0\in \sigma(B)$  for which the quasi-nilpotent part of $B$ vanishes. Then $ \exists C \geq 0$ so that for any $z \neq 0$ we have for sufficienly large $\beta \gg 1$ that $\beta B - z$ is boundedly invertible. Furthermore we have  the following  convergence  towards the Riesz projector $P = P_{\{0\}} $:
	\begin{align}\label{convergence_to_projection}
		\left\|(\beta B - z)^{-1} - P/(-z) \right\| \leq C/\beta \longrightarrow 0.
	\end{align}
	If $B P \neq 0$, we have $\|(\beta B - z)^{-1} \| \rightarrow \infty$ instead.
\end{Lem}

\begin{proof}
	First, by assumption, since $0 \in \sigma(B)$ is isolated, 
	there exists some $\epsilon > 0$ such that the punctured disk
	$
	\mathbb{D}_\epsilon \setminus \{0\} = \{ \lambda \in \mathbb{C} : 0 < |\lambda| < \epsilon \}
	$
	lies entirely within the resolvent set $\rho(B)$ of $B$. Bounded invertibility of $\beta B - z$ is equivalent to bounded invertibility of
	$
	B - z/\beta
	$
	and for $|z| \leq \beta \epsilon$, $ B - z/\beta$ is boundedly invertible.
	
	
	
	To prove the second claim, we leverage the spectral decomposition induced by the Riesz projection $P$:
	On its domain $\mathcal{D}(B)$, we can decompose $B$ into a sum of closed operators as
	$
	B = P B P + Q B Q,
	$
	where $Q = I - P$ is the complementary projection. 
	Since $PQ = QP = 0$, this decomposition induces a corresponding block decomposition of the resolvent as
	$
	(\beta B - z)^{-1} = P (\beta P B P - z)^{-1} P + Q (\beta Q B Q - z)^{-1} Q.
	$
	By assumption, the quasi-nilpotent part of $B$ vanishes, so $P B P = 0$. Hence,
	\[
	P (\beta P B P - z)^{-1} P = P (-z)^{-1} P = - P/z.
	\]
	
	Therefore
	$
	-z (\beta B - z)^{-1} = P + (-z) Q (\beta Q B Q - z)^{-1} Q.
	$	
	To conclude the proof, it suffices to show that there exists $C \geq 0$ such that
	$
	\left\| Q (\beta Q B Q - z)^{-1} Q \right\| \leq C/\beta
	$
	for sufficiently large $\beta$.
	To verify this estimate, note that
	$
	Q (\beta Q B Q - z)^{-1} Q = \frac{1}{\beta} Q \left( Q B Q - \frac{z}{\beta} \right)^{-1} Q.
	$
	Since $0$ is isolated in $\sigma(B)$, and $P$ projects onto the eigenspace of $0$, the operator $Q B Q$ has a resolvent set containing a neighborhood of $0$. In particular, there exists $\epsilon > 0$ such that the closed disk of radius $\epsilon$ centered at $0$ is contained in the resolvent set of $Q B Q$.
	For $|\frac{z}{\beta}| < \epsilon$ we can thus -- by analyticity of the resolvent -- expand it in a Neumann series:
	\[
	Q \left( Q B Q - \frac{z}{\beta} \right)^{-1} Q = \sum_{n=0}^\infty \left( \frac{z}{\beta} \right)^n \left[ Q \left( Q B Q \right)^{-1} \right]^{n+1} Q,
	\]
	Thus,
	%		\begin{align}
		$
		\left\| Q (\beta Q B Q - z)^{-1} Q \right\| 
		\leq \frac{1}{\beta} \left\| Q (Q B Q)^{-1} \right\| \sum_{n=0}^\infty \left| z/\beta \right|^n \left\| Q (Q B Q)^{-1} \right\|^n$.
	For sufficiently large $\beta$, the geometric series converges and is bounded by $2$.
	Thus we find
	$
	\left\| Q (\beta Q B Q - z)^{-1} Q \right\| \leq 2 \left\| Q (Q B Q)^{-1} \right\|/\beta.
	$
	Setting
	$
	C := 2 \left\| Q (Q B Q)^{-1} \right\|
	$
	completes the proof.
	
	
	
	It remains to prove that whenever $B P \neq 0$, we have $\|(\beta B - z)^{-1} \| \rightarrow \infty$ instead. We note $\|(\beta B - z)^{-1} \| \geq [ \|(\beta B - z)^{-1}P \| /\|P\|= \|(\beta BP - z)^{-1}P \|/\|P\| $. Suppose $BP \neq 0$ and take $\psi \in \mathcal{D}(B)$ so that $\|P\psi\| = 1$ and $BP\psi \neq 0$. Set $\eta_{\beta,z} = (\beta BP -z)^{-1} P\psi$. Then $B P\psi = (\eta_{\beta,z} + z \psi)/\beta$. Thus we find $\|\eta_{\beta,z} \|/\beta \rightarrow  \|B P\psi\| \neq 0$ as $\beta \rightarrow \infty$ and hence $\|\eta_{\beta,z}\| \rightarrow \infty$. Thus, we have $\|(\beta B - z)^{-1} \| \gtrsim \|(\beta BP - z)^{-1}P \| \geq \|\eta_{\beta,z}\| \rightarrow \infty$.
	
\end{proof}

\end{document}
